\section{Removing the determinant hypothesis from an odd adjoint origin}
\label{sec:kt-odd-adjoint-scalar-normalization}

The determinant normalization in
KT-ADJOINT-BRANCH-UNIT-EXCLUSION can be recovered by a global scalar
twist.  The input needed here is an odd rank-two origin over
$\mathbf Q$.  Descent of a self-dual automorphic rank-three system to
such an origin is a separate assertion.

\begin{theorem}[KT-ODD-ADJOINT-SCALAR-NORMALIZATION]
Let $p>3$, let $D<0$ be a fundamental discriminant with $p\nmid D$,
put $E=\mathbf Q(\sqrt D)$ and $\eta=\chi_D$, and let $e\ge2$ be
even with $(p-1)\nmid2e$.  Let $V$ be a continuous absolutely
irreducible two-dimensional representation of $G_{\mathbf Q}$ over
a finite extension of $\mathbf Q_p$.  Suppose $V$ is projectively
odd: complex conjugation acts with two distinct eigenvalues.
Put $A=\operatorname{Ad}^0(V)$.  Assume that
\begin{enumerate}
\item $A|_{G_E}$ is unramified at every finite place away from $p$;
\item $(\bar A|_{G_E})^{\mathrm{ss}}
       =1\oplus\omega^e\oplus\omega^{-e}$;
\item at a place $v\mid p$ of $E$, $A|_{G_{E_v}}$ has a full stable
      flag with ordered inertial characters
      $(\varepsilon_p^e,1,\varepsilon_p^{-e})$.
\end{enumerate}
Then, after a finite coefficient extension, there is a continuous
character $\theta:G_{\mathbf Q}\to\overline{\mathbf Q}_p^{\times}$,
with image in a finite coefficient field, such that
$V'=V\otimes\theta^{-1}$ satisfies
\[
 \det V'=\eta\varepsilon_p^e,\qquad
 (\bar V')^{\mathrm{ss}}=1\oplus\bar\eta\omega^e.       \tag{SN.1}
\]
Moreover, $V'|_{G_E}$ is unramified away from $p$, and $V'$ has a
$G_{\mathbf Q_p}$-stable line with unramified rank-one quotient.
Consequently $B_{1,\eta\omega^e}$ is not a $p$-adic unit.
There is no upper bound on $e$, and $p$ may split or be inert in $E$.
\end{theorem}

\begin{proof}
All lattices, eigenlines and characters below are defined after
finite coefficient extensions.  Their finiteness, including that
of the eventual global character, is justified in the construction.
Write $\psi=\omega^e$.
Since $p\nmid D$, $E\cap\mathbf Q(\mu_p)=\mathbf Q$ and the order
of $\psi|_{G_E}$ is greater than two.

\paragraph{The residual scalar character and its orientation.}
The residual-projective argument in the proof of
KT-ADJOINT-BRANCH-UNIT-EXCLUSION, up to (AJ.2), uses only
hypothesis~2 and this order condition.  It gives
\[
 \bar V^{\mathrm{ss}}=\lambda\oplus\mu,
 \qquad \lambda\mu^{-1}=\bar\eta^t\psi,
 \quad t\in\{0,1\}.                                    \tag{SN.2}
\]
For clarity, that argument first triangularizes the residual
adjoint on $G_E$.  Its faithful projective image is a finite group
with a normal unipotent $p$-subgroup and abelian prime-to-$p$
quotient.  A nontrivial such $p$-subgroup has a unique fixed point
on $\mathbf P^1$; if it is trivial, an element of order greater
than two forces the projective image into a torus.  Thus the two
residual characters exist on $G_E$.  Their ratio is conjugation
invariant and different from its inverse, so the global action
cannot exchange them.  This gives (SN.2).

Because complex conjugation is an involution and is not scalar
on $V$, its eigenvalues are $1,-1$.  Reduction preserves their
distinctness as $p>2$.  Since $e$ is even, $\psi(c)=1$ and
$\eta(c)=-1$, so (SN.2) has $t=1$.  Hence
\[
          \bar V^{\mathrm{ss}}=\mu\bar\eta\psi\oplus\mu.
                                                               \tag{SN.3}
\]
Here $\mu$ is an arbitrary finite residual character; it is not
assumed quadratic.

\paragraph{The quotient character at $p$.}
The highest line of the flag in hypothesis~3 is the unique inertia
eigenline of $A$ with character $\varepsilon_p^e$.
It is stable under $G_{\mathbf Q_p}$, since $E_v/\mathbf Q_p$ is
unramified or trivial and the cyclotomic character is invariant
under conjugation.  A nonzero element $N$ of that line is a
trace-zero matrix.  Conjugation preserves its determinant,
whereas inertia multiplies $N$ by the infinite-order character
$\varepsilon_p^e$.  Therefore $\det N=0$.
The line $\ker N=\operatorname{im}N$ in $V$ is
$G_{\mathbf Q_p}$-stable.  If $\delta_+$ and $\delta_-$ denote
its character and the quotient character, respectively, then
\[
            \delta_+/\delta_-\big|_{I_p}
                        =\varepsilon_p^e\big|_{I_p}.     \tag{SN.4}
\]
Intersecting this line with a stable lattice and taking the
saturated quotient identifies its reduced characters with the
two factors in (SN.3) on inertia.  Since $\eta$ is unramified
at $p$ and $\psi^2|_{I_p}\ne1$, the order in (SN.4) forces
\[
                   \bar\delta_-|_{I_p}=\mu|_{I_p}.        \tag{SN.5}
\]
Indeed the opposite order would give ratio $\psi^{-1}$.
We will prescribe $\theta|_{I_p}=\delta_-|_{I_p}$.

\paragraph{Compatible scalar choices at every other prime.}
Fix a rational prime $q\ne p$, and choose a decomposition group
$G_{\mathbf Q_q}$ and its inertia group $I_q$.
For $g\in I_q\cap G_E$, hypothesis~1 and the kernel of the
adjoint action give $V(g)=z_g I$.

If $q\nmid D$, this says that all of $V(I_q)$ is scalar.
A Frobenius lift has an eigenline after a finite coefficient
extension.  That line is stable under inertia and the Frobenius
lift, hence under the full decomposition group, whose quotient
by inertia is topologically procyclic.  Its character
$\alpha_q$ therefore satisfies
\[
               \alpha_q(g)=z_g\quad(g\in I_q),
               \qquad \bar\alpha_q|_{I_q}=\mu|_{I_q}.     \tag{SN.6}
\]

If $q\mid D$, the projective inertia image has order at most two.
It has order exactly two: (SN.3) has distinct inertial characters
$\mu\bar\eta$ and $\mu$, since $q\ne p$.
Choose $h\in I_q$ with $\eta(h)=-1$.  The matrix $V(h)$ has two
distinct eigenvalues with ratio $-1$.  They have distinct
reductions, so their spectral idempotents are integral on the
chosen stable lattice after a finite extension.  Every inertia
matrix is scalar or a scalar multiple of $V(h)$; thus there are
two inertia eigenlines, one with reduced character $\mu|_{I_q}$
and the other with $\mu\bar\eta|_{I_q}$.
The decomposition group preserves each line.  To see this,
conjugation by any of its elements sends the first character to
one whose reduction is still $\mu|_{I_q}$, because $\mu$ is a
global character.  The second line has a distinct reduction, so
the two lines cannot be exchanged.
Let $\alpha_q$ be the character of the first line.  Then
\[
 \bar\alpha_q|_{I_q}=\mu|_{I_q},\qquad
 \alpha_q(g)=z_g\quad(g\in I_q\cap G_E),                 \tag{SN.7}
\]
and the ratio of the other inertia character to $\alpha_q$
is exactly $\eta|_{I_q}$.  It is trivial on $I_q\cap G_E$
and equals $-1$ on the other coset.

\paragraph{Only finitely many of these inertia characters are nontrivial.}
We give this step explicitly; continuity alone is not being used
as a general finite-ramification assertion for arbitrary
$p$-adic representations.  For $q\ne p$, the restriction of a
continuous character of $G_{\mathbf Q_q}$ to inertia has finite
image in any finite extension of $\mathbf Q_p$.
Local class field theory identifies that image as a quotient of
$\mathbf Z_q^{\times}$.  An open pro-$q$ subgroup has trivial
image in a sufficiently small torsion-free pro-$p$ subgroup of
the target; the remaining quotient is finite.
In particular, when $q\nmid Dp$, the scalar image $V(I_q)$
consists of roots of unity in the original coefficient field $K$.

Choose $n$ such that $1+\varpi^n\mathcal O_K$ is torsion free.
The finite representation $V\bmod\varpi^n$ is unramified
outside a finite set of rational primes.  Outside that set and
the divisors of $Dp$, the scalar image $V(I_q)$ is both torsion
and contained in $1+\varpi^n\mathcal O_K$, so is trivial.
Together with (SN.6)--(SN.7), this proves that
$\alpha_q|_{I_q}=1$ for all but finitely many $q\ne p$.

\paragraph{A single global character with the prescribed inertias.}
For each relevant prime $q$, local class field theory for
$\mathbf Q_q$ says that the inertia restriction of a local
character factors continuously through the $q$-adic cyclotomic
character
\[
                 \varepsilon_q:I_q\longrightarrow
                                      \mathbf Z_q^{\times}.
\]
Indeed the maximal abelian extension is the compositum of the
unramified extension and $\mathbf Q_q(\mu_{q^{\infty}})$;
the latter is totally ramified.  Thus there is a continuous
$\beta_q:\mathbf Z_q^{\times}\to K'^{\times}$ such that
\[
 \beta_p\circ\varepsilon_p|_{I_p}=\delta_-|_{I_p},
 \qquad
 \beta_q\circ\varepsilon_q|_{I_q}=\alpha_q|_{I_q} (q\ne p).
                                                               \tag{SN.8}
\]
The field $K'$ can be chosen finite: only finitely many nontrivial
characters occur, and each was obtained from an eigenline over
a finite extension.  Set
\[
               \theta=\prod_q\bigl(\beta_q\circ
                         \varepsilon_q:G_{\mathbf Q}\to K'^{\times}\bigr),
                                                               \tag{SN.9}
\]
omitting the trivial factors.  This is a finite product of
continuous characters.  At $I_q$ every factor with a different
index is unramified, so (SN.9) has exactly the inertias prescribed
in (SN.8).  There is no separate choice at the real place.

Equations (SN.5)--(SN.7) show that $\bar\theta\mu^{-1}$ is
unramified at every finite prime of $\mathbf Q$.
Such a finite character is trivial: its fixed field would be
a finite extension of $\mathbf Q$ unramified at every finite
prime, which is trivial by Minkowski's discriminant bound.
Consequently $\bar\theta=\mu$, giving the residual statement
in (SN.1).

\paragraph{The determinant and the ordinary normalization.}
At $p$, twisting by (SN.9) makes the quotient character
unramified; (SN.4) makes the determinant on inertia equal to
$\varepsilon_p^e$.  At $q\nmid Dp$, all inertia matrices of
$V'$ are the identity.  At $q\mid D$, their two characters are
$1,\eta|_{I_q}$ by (SN.7).  It follows that
\[
                    \det V'\,(\eta\varepsilon_p^e)^{-1}
                                                               \tag{SN.10}
\]
is a continuous character unramified at all finite primes.
It too is trivial: every finite quotient of its compact image
would give a finite everywhere-unramified extension of
$\mathbf Q$.  This proves the determinant statement of (SN.1).
Equations (SN.6)--(SN.8) also show directly that
$V'|_{G_E}$ is unramified away from $p$.

Absolute irreducibility is unchanged by the scalar twist.
The residual representation, ordinary quotient and away-prime
unramifiedness now meet all hypotheses of
KT-ORDINARY-RANK2-UNIT-EXCLUSION with $a=e$.
That theorem gives the asserted nonunit conclusion.
\end{proof}

\paragraph{Source scope and imported class field theory.}
The local character factorization used in (SN.8) is the
$\mathbf Q_q$ specialization of J.~S.~Milne,
\emph{Class Field Theory}, version~4.03 (6 August 2020),
Chapter~I, \S1.10--1.11, printed pp.~23--24, especially
the displayed decomposition of $\mathbf Q_q^{\mathrm{ab}}$.
The local Kronecker--Weber statement is Theorem~4.8 and
Corollary~4.12, pp.~47--48.  The everywhere-unramified argument
is the one in Lemma~4.17 and its proof, pp.~49--50.
These are imported class field theory and discriminant inputs;
no new proof of them or Lean kernel check is claimed here.
The author PDF is preserved in the adjacent primary-source
directory with its retrieval hash.

An actual Type~I system with an identified odd $G_{\mathbf Q}$
adjoint origin therefore needs no additional hypothesis on the
origin's determinant, nebentypus or scalar twist.  This theorem
does not by itself construct such an origin from a
conjugate-self-dual rank-three representation over $E$.
